\DocumentMetadata{lang=en-US,pdfstandard=ua-2,tagging=on}
\documentclass{ximera}
\usepackage{unicode-math}
\usepackage[points]{ximeraExam}


%%%%%%%%%%%%
%%
%% This is an example demonstrating features of the package, not a minimal starter template.
%%
%% Demonstrates:
%%   - PDF/UA-2 document setup
%%   - custom exam cover and grade table
%%   - inherited question points
%%   - nested parts
%%   - horizontal multiple choice
%%   - answer lines and student response space
%%   - running headers/footers
%%   - TikZ/pgfplots content
%%
%%%%%%%%%%%%



% preliminary setup for ximeraExam
\noprintanswers
\renewcommand{\pointsLabel}[1]{(#1 points)}
\answerlineAlign{right}


% Update the course, date, and assignment time limit here
\newcommand*{\courseName}{Calculus 1} 
\newcommand*{\monthAndDate}{February 17}
\newcommand*{\timeLimit}{45 Minutes}




\newcommand*{\courseSemester}{Spring 2026}
\newcommand*{\examName}{Midterm 2}
\newcommand*{\formDesignation}{Form B}
\newcommand*{\dueDate}{\monthAndDate , 2026 }

%set the title
\title{\courseName \,\examName \, \courseSemester}


% data used by ximeraExam to set pdf title
% also used if sticking with default examcover
\examtitle{\examName}
\examcourse{\courseName - \courseSemester}
\examdate{\dueDate}
\examversion{\formDesignation}


% question/part labels and set headers
\renewcommand{\questionlabel}{\textbf{\questionnumber)}}
\renewcommand{\partlabel}{\partnumber)}
\firstpageheader{}{}{}
\runningheader{\examName}{\courseName \, \, \courseSemester}{\formDesignation}
\runningfooter{}{}{}





% additional packages
\usepackage{pgfplots}   % <- for graphics
\pgfplotsset{compat=newest}
\AddToHook{enddocument/afterlastpage}{\label{LastPage}}
\usetikzlibrary{shapes.geometric, calc}
\usepackage[left=0.5in, right=0.5in, bottom=0.5in, top=0.7in, headsep=1.5em]{geometry}

% some notation for questions
\newcommand{\ddx}{\frac{d}{d x}} 
  
% For points in graphs
\newcommand{\opencircle}[1]{	
 						\draw[fill=blue] #1 circle [color=blue,radius=3pt];
						\draw[fill=white] #1 circle [color=white,radius=2pt]; }
						
\newcommand{\closedcircle}[1]{	
 						\draw[fill=blue] #1 circle [color=blue,radius=3pt];
						}
 
\newcommand{\noWork}{(You do not need to show work for this question.)}





\begin{document}
  % Set sans serif fonts for readability.
	\sffamily


\begin{examcover} %using a custom examcover page

	% Modify styling of cover page.
	\thispagestyle{empty}
	\newgeometry{margin=0.5in}
	 
	
	% Start of Cover Page
	
	\examcoverheading{\textbf{\courseName \, \examName}}
	
	% name info as text instead of a tabular, so it's not mistaken as a data table.
	\begin{flushleft}
	\noindent
	\makebox[0.30\linewidth][l]{\textbf{\formDesignation}}%
	\hfill Printed First and Last Name: \rule{2.5in}{.7pt}\par
	\medskip
	
	\noindent
	\makebox[0.30\linewidth][l]{\dueDate}%
	\hfill OSU name.\#: \rule{2.5in}{.7pt}\par
	\medskip
	
	\noindent
	\makebox[0.30\linewidth][l]{Duration: \timeLimit}%
	\hfill Lecturer: \rule{2.5in}{.7pt}\par
	\medskip
	
	\noindent
	\makebox[0.30\linewidth][l]{Page 1 of \pageref{LastPage}}%
	\hfill Recitation Instructor: \rule{2.5in}{.7pt}\par
	\medskip
	
	\noindent
	\hfill Recitation Time: \rule{2.5in}{.7pt}\par
	\end{flushleft}
	
	
	\vspace*{\stretch{1}}
	 
	\noindent \textbf{Instructions:} 
	\begin{itemize}
		\item These questions are here for you to demonstrate your understanding of the course material.\\
	 		To do that, \textbf{show all relevant supporting work} to receive full credit on all Problems unless explicitly stated not to.
	 		
		\item Incorrect answers with substantially correct work may receive partial credit.  \\
			\textbf{Answers without justification will receive no credit.}\\
			Using a table of values or L'H\^{o}pitals Rule are \textbf{not justification} for computing limits.
		
		\item Give \textbf{exact answers} unless instructed to do otherwise.\\
			Intervals should be expressed \textbf{simplified in interval notation}.\\
			If a requested value does not exist, write ``DNE".
			Write that limits ``DNE" only if they do not exist and are neither $\infty$ nor $-\infty$.
	
		\item You are expected to \textbf{evaluate any famous functions at standard values}.
		
		\item This is a closed book/closed note exam. \\
			Scrap paper is not allowed. All your work should be on this exam paper.\\
			Calculators, tablets, headphones, and other devices are not permitted.
		
	
	

	\end{itemize}
		
	\vspace*{\stretch{1}}
	
	% Box for students to sign acknowledging the instructions.
	\begin{tcolorbox}[colback=white]
		\textbf{SIGN YOUR NAME:}  \mbox{\begin{minipage}{50ex}\hfill\vspace{3em}\end{minipage}}\\
	\end{tcolorbox}
	
	\vspace*{\stretch{1}}
	
	\begin{center}
		% setup the gradetable. Horizontal, and not showing the ``Earned'' row, since not used here for hand scoring.
		\gradetable[h][noearned]
	\end{center}
	
	\vspace*{\stretch{1}} 
	

\end{examcover}

%\newpage
\begin{questions}
	

\newcommand{\ptXMin}{-4}
\newcommand{\ptXMax}{5}
\newcommand{\ptYMin}{-3}
\newcommand{\ptYMax}{4}
\newcommand{\DomF}{(\ptXMin, \ptXMax)}
\newcommand{\secA}{-3}
\newcommand{\secB}{-1}

\newcommand{\CompPtA}{\ensuremath{-3.8}}
\newcommand{\CompPtB}{\ensuremath{-1.7}}
\newcommand{\CompPtC}{\ensuremath{4.3}}
%Question 1
\question 
The entire graph of a function $f$, with domain $\DomF$, is given below.



%% Tikz graph with alt-text.
%% Give the complete TikZ graphic meaningful alt text for tagged PDF output.
%% Describe the mathematically relevant content rather than only saying
%% "graph of f."

\begin{center}
	\begin{tikzpicture}[
    alt={
    Graph of a function f on the domain (-4,5), with several
    connected curve segments and open or closed endpoints.
  }
]
		\begin{axis}[
			xmin=\ptXMin.3, xmax=\ptXMax.3, ymin=\ptYMin.2,ymax=\ptYMax.2,    
			axis lines =middle, 
			every axis y label/.style={at=(current axis.above origin),anchor=south},
			every axis x label/.style={at=(current axis.right of origin),anchor=west},
			xtick={\ptXMin,...,\ptXMax}, ytick={\ptYMin,...,\ptYMax},
			grid=major, width= 0.8*\pagewidth, height = 3.5in,
			grid style={dashed, gray!40}
			]
			
			\addplot[color=blue, ultra thick, smooth, domain=-4:-2]{2*(x+4)-1};
			\addplot[color=blue, ultra thick, smooth, domain=-2:0]{(x^2-1)};
			\addplot[color=blue, ultra thick, smooth, samples=200, domain=0:1]{-(1-x^2)^(0.5)};
			\addplot[color=blue, ultra thick, smooth, samples=500, domain=1:5]{(x-1)^(0.5) };

			\opencircle{(-4,-1)};
			\closedcircle{(-1,-2)};
			\opencircle{(-1, 0)};
			\opencircle{(5,2)};				
		\end{axis}
	\end{tikzpicture}
\end{center}

\begin{parts}
	\part[3] At which point(s) in $\DomF$ is $f$ \textbf{nondifferentiable}? \noWork
	
	\begin{solution}[\stretch{1}]
	\end{solution}	
	\answerline[width=2in]{At $x=$}
	
	\part[3] \textbf{Sketch} the line \textbf{secant} to the graph of $f$ at the points corresponding to $x=\secA$ and $x=\secB$ on the axes above.  \textbf{Label it as $S$.}	\noWork	

	\begin{solution}[\stretch{1}]
	\end{solution}	
	
		
	\part[2] Which one of these values is \textbf{least}? \noWork 

	\vspace{0.25in}
	
	\begin{multipleChoice}[layout=horizontal, max-per-row=3]
	    \choice{ \(f'(\CompPtA)\) }
	    \choice{ \(f'(\CompPtB)\) }
	    \choice{ \(f'(\CompPtC)\) }
	\end{multipleChoice}

	
	\begin{solution}[\stretch{1}]
	\end{solution}	

	\answerline[width=1in]{Answer (A, B, or C):}
	
	
	\part[2] Which one of these values is \textbf{greatest}?\noWork
	
	\vspace{0.25in}
	
	\begin{multipleChoice}[layout=horizontal]
	    \choice{ \(f'(\CompPtA)\) }
	    \choice{ \(f'(\CompPtB)\) }
	    \choice{ \(f'(\CompPtC)\) }
	\end{multipleChoice}
	
	\begin{solution}[\stretch{1}]
	\end{solution}	
	
	\answerline[width=1in]{Answer (A, B, or C):}
	
\end{parts}

\newpage

\newcommand*{\DoDa}{6}
\newcommand*{\fnPos}{\ensuremath{2\sqrt{3+t}}}
\question Suppose an object is traveling along a line with position function given by $s(t) = \fnPos$ meters, for $t \geq 0$ in hours.

\vspace{1em}

\begin{parts}
	\part[2] Evaluate $s(\DoDa)$.
	
	\begin{solution}[\stretch{1}]
	\end{solution}	
	\answerline[width=2in]{Value:}

	
	\part[3] For this function, calculate $s(\DoDa+h)$.
	
	\begin{solution}[\stretch{1}]
	\end{solution}	
	\answerline[width=2in]{Result:}

	
	\part[7]\label{PART:limit} Find the form and calculate the limit: $\displaystyle \lim_{h \to 0} \dfrac{s(\DoDa+h)-s(\DoDa)}{h}$.			
	
	\begin{solution}[\stretch{5}]
	\end{solution}	
	\answerline[width=2in]{Form:}
	
	\vspace{0.25in}

	\answerline[width=2in]{Value:}


	\part[4] Explain what the value you calculated in \ref{PART:limit}) represents about the object's motion. Be specific.
	\begin{solution}[\stretch{2}]
	\end{solution}	 
	
	
	
	
\end{parts}

				


\newpage


\newcommand{\tlFn}{g}
\newcommand*{\fnOrig}{\ensuremath{\dfrac{3x+2}{4x+1}}}
\newcommand*{\fnDeriv}{\ensuremath{ -\dfrac{5}{(4x+1)^2}}}

\newcommand{\pointA}{2}
\newcommand{\pointB}{\dfrac{8}{9}}
\newcommand{\toPt}{1}
\newcommand{\innerFn}{3x^3-1} % needs to have value \pointA when x is \toPt

\question Let $\tlFn$ be the function given by $\tlFn(x) = \fnOrig$.

\vspace*{0.15in}
\begin{parts}
	
	\part[5] Show that the derivative of $\tlFn$ is given by $\fnDeriv$

	(HINT: Use Derivative shortcut formulas.)
	
	\begin{solution}[\stretch{2}]
	\end{solution}		
	
	\part[4] Find an equation for the line \textbf{tangent} to the graph of $y=\tlFn(x)$ at the point $\left( \pointA, \pointB \right)$. 

	\begin{solution}[\stretch{1}]
	\end{solution}	
	
		\answerline[width=4.5in]{Equation:}


	
	\part[6] Evaluate the limit: $\displaystyle \lim_{x \to \toPt} \dfrac{\tlFn(\innerFn)-\pointB}{x-\toPt}$ . \textbf{EXPLAIN}.

	\begin{solution}[\stretch{1}]
	\end{solution}	
	\answerline[width=4.5in]{Value:}
	
	
\end{parts}		


\newpage
\newcommand{\tabFn}{\ensuremath{k}}
\newcommand{\tabX}{\ensuremath{3}}
\newcommand{\tabY}{\ensuremath{4}}


\question A function $\tabFn$, with domain $(-\infty, \infty)$, has values and derivative values given in the table below.
\[ \begin{array}{|c|c|c|c|c|c|c|c|c|c|}
	\hline
	\mathbf{x}    & -1&  0 & 1 & 2 & 3 & 4 & 5 & 6 & 7   \\ \hline
	\mathbf{\tabFn(x)} & 5 & 8 & -3  & 0  &4& 6 & 7 & -\sqrt{3} & -2 \\ \hline
	\mathbf{\tabFn'(x)}& 2 & -4 & 1 & \frac{2}{3}  &-5 & \frac{1}{4} & -2 &  -6 & 1 \\ \hline
\end{array} \]

\begin{parts}
	\part[4] Using the provided information, we can conclude that $\displaystyle \lim_{x\to \tabX} \tabFn(x) = \tabY$. \textbf{EXPLAIN} how we know this is true.
	
	\begin{solution}[\stretch{1}]
	\end{solution}	
		
		
	\part[7] Calculate $\displaystyle \left[ \ddx\left( \tabFn(x)  e^{3x} - 6x \right) \right]_{x=5}$
	
	\begin{solution}[\stretch{1}]
	\end{solution}	
	
	\answerline[width=2.5in]{Value:}

	
	
	\part[8] Calculate $\displaystyle \left[ \ddx \left( (x^2+1)^{12} + \tabFn\left( 2\sin(x) + 5 \right)  \right)  \right]_{x=0}$
	
	\begin{solution}[\stretch{1}]
	\end{solution}	

	\answerline[width=2.5in]{Value:}


	
	%	
\end{parts}	




\end{questions}

\end{document}
